\documentclass[a4paper,12pt]{article}
\usepackage{graphicx}
\usepackage{hyperref}
\usepackage{geometry}
\usepackage[ngerman]{babel}
\geometry{margin=1in}

% Title Page
\title{\textbf{Anleitung zum Raspi-CyberSec-Lab}}
\author{Jonas Schmitt}
\date{\today}

\begin{document}

\maketitle

\tableofcontents
\newpage

\section{Allgemeine Informationen}
Das RaspberryPI Cybersecurity Lab ist eine Platform zum Testen verschiedener Cyberangriffe, das Projekt ist Teil meiner Bachelorarbeit in Elektro- und Informationstechnik an der Technischen Hochschule Nürnberg.
Es hat verschiedene Funktionen um Umgebungen zu schaffen, welche für das Ausprobieren von Angriffen aus den Bereichen Wifi Netzwerke, Webseiten / -applikationen, MQTT Kommunikation und Bluetooth ausgelegt sind.

\section{Voraussetzungen}
Das Projekt wurde mit folgender Hardware umgesetzt:
\begin{itemize}
    \item RaspberryPI 4B mit 8GB RAM
    \item Fenvi AX1800 USB Netzwerkkarte
    \item Waveshare ESP32c6 Microcontroller
    \item Waveshare 4.3 zoll LCD 
    \item Drehgeber mit Druckknopf
\end{itemize}

Als Betriebssystem des RaspberryPI dient RaspberryPI OS in der light Version (nur Kommandozeile).
Zum Kompilieren und Flashen des ESP32 Codes wird das ESP-IDF benötigt, dieses kann in Visual Studio Code als Erweiterung installiert werden.

\section{Installation}
 Zuerst muss das Github \href{https://github.com/Der-Erzfeind/Raspi-CyberSec-Lab-Project.git}{Repository} des Projekts auf den RaspberryPI geklont werden, dies kann mit folgendem Kommando erfolgen:
 \begin{verbatim}
    git clone https://github.com/Der-Erzfeind/Raspi-CyberSec-Lab-Project.git
 \end{verbatim}
Es wird empfohlen das Repository in das standardmäßige Nutzerverzeichnis zu klonen.

Nach dem Herunterladen der Projektdaten muss das Installationsskript setup.sh ausgeführt werden.

\section{Benutzeranleitung}

Zum Einschalten des Geräts, schließen Sie das Gerät mit dem zugehörigen Netzteil an den Strom an.
Das Gerät fährt nun hoch, wobei der Bootvorgangn zu sehen ist. (in der linken oberen Ecke wird das RaspberryPI Logo angezeigt)

Zur Navigation wird das Drehrad genutzt:
\begin{figure}
    \item 
\end{figure}

\subsection{Übersicht Hauptmenü}
Nach erfolgreichem Bootvorgang wird das Hauptmenü angezeigt, in diesem stehen folgende Optionen zur Verfügung:
\begin{itemize}
    \item \textbf{Wifi} - Optionen zu Wifi Netzwerken
    \item \textbf{Bluetooth} - Optionen zu Bluetooth
    \item \textbf{Webapp} - Optionen zu Webapplikationen
    \item \textbf{power off} - Ausschalten des Geräts
\end{itemize}

\section{Features and Usage}
\subsection{Feature 1}
Description and steps to use Feature 1.

\subsection{Feature 2}
Description and steps to use Feature 2.

\section{Troubleshooting}
Common issues and their solutions.
\begin{itemize}
    \item \textbf{Issue 1:} Description and solution.
    \item \textbf{Issue 2:} Description and solution.
\end{itemize}

\section{Support and Contact Information}
For further support, contact:
\begin{itemize}
    \item Email: support@example.com
    \item Phone: +1-234-567-890
    \item Website: \url{http://example.com/support}
\end{itemize}

\end{document}
